\VerdierPhysicalPageAnchor{0077}%
\VerdierSemanticAnchor{ch1-2-4-3-simple-complex-morphisms}
D'où, en passant aux complexes simples, deux morphismes de foncteurs :

\VerdierSemanticAnchor{ch1-2-4-3-eq2}
\begin{equation}
  \begin{aligned}
    \operatorname{comp}_{\Pi}h\colon
      \operatorname{comp}_{\Pi}F
      &\longrightarrow \operatorname{comp}_{\Pi}F',\\
    \operatorname{comp}_{\Sigma}h\colon
      \operatorname{comp}_{\Sigma}F
      &\longrightarrow \operatorname{comp}_{\Sigma}F'.
  \end{aligned}
  \tag{2.4.3.2}
\end{equation}

\VerdierSemanticAnchor{ch1-2-4-3-composition}
Si $h$ et $h'$ sont deux morphismes composables de foncteurs multiadditifs,
on a :
\begin{equation*}
  \begin{aligned}
    \operatorname{comp}_{\Pi}hh'
      &=\operatorname{comp}_{\Pi}h\,\operatorname{comp}_{\Pi}h',\\
    \operatorname{comp}_{\Sigma}hh'
      &=\operatorname{comp}_{\Sigma}h\,\operatorname{comp}_{\Sigma}h'.
  \end{aligned}
\end{equation*}

\VerdierSemanticAnchor{ch1-2-4-3-sign-rule-ii}
Le morphisme $\operatorname{comp}_{\Pi}h$ (resp.
$\operatorname{comp}_{\Sigma}h$) est un morphisme de
$\mathbb Z[n]$-foncteurs qui suit la règle des signes II (1.6.6).

\VerdierSemanticAnchor{ch1-2-4-4}
\noindent\textbf{2.4.4.}\quad
Lorsque le foncteur $F$ est contravariant en certaines variables,
\emph{i.e.} lorsqu'on a des égalités
$C_i=(C_i')^\circ$, pour $i\in J\subset[n]$, les foncteurs
$\operatorname{comp}_{\Pi}F$ et $\operatorname{comp}_{\Sigma}F$ sont,
modulo les identifications faites en (2.4.1), contravariants en les variables
correspondantes.

\VerdierSemanticAnchor{ch1-2-5}
\subsection{Homotopies}

\VerdierSemanticAnchor{ch1-2-5-1}
\noindent\textbf{2.5.1.}\quad
Soient $X$ et $Y$ deux complexes $n$-uples de $C$ et
$f,g\colon X\rightrightarrows Y$ deux morphismes de complexes. Une
\emph{homotopie} reliant $f$ à $g$ est une famille de morphismes

\VerdierSemanticAnchor{ch1-2-5-1-family}
\begin{equation*}
  a=\bigl(a^{\sigma,i}\colon X^{\sigma+e_i}\longrightarrow Y^\sigma,
    \ \sigma\in\mathbb Z[n],\ i\in[n]\bigr)
\end{equation*}
telle que :

\VerdierSemanticAnchor{ch1-2-5-1-eq1}
\begin{equation}
  \begin{aligned}
    a^{\sigma+e_i-e_j,j}d_X^{\sigma,i}
      &=d_Y^{\sigma-e_j,i}a^{\sigma-e_j,j},
      &&\forall\sigma\in\mathbb Z[n],\
        \ \forall i,j\in[n],\ i\ne j,\\
    g^\sigma-f^\sigma
      &=\sum_{i\in[n]}\Bigl(
          a^{\sigma,i}d_X^{\sigma,i}
          +d_Y^{\sigma-e_i,i}a^{\sigma-e_i,i}\Bigr),
      &&\forall\sigma\in\mathbb Z[n].
  \end{aligned}
  \tag{2.5.1.1}
\end{equation}

\VerdierSemanticAnchor{ch1-2-5-1-notation}
On utilise la notation $a\colon f\rightsquigarrow g$ pour désigner les
homotopies $a$ reliant $f$ à $g$.

\VerdierSemanticAnchor{ch1-2-5-2}
\noindent\textbf{2.5.2.}\quad
Soient $f,g,h\colon X\longrightarrow Y$ trois morphismes,
$a\colon f\rightsquigarrow g$ et $b\colon g\rightsquigarrow h$ deux
homotopies. On désigne par $b+a\colon f\rightsquigarrow h$ l'homotopie

\VerdierSemanticAnchor{ch1-2-5-2-sum}
\begin{equation*}
  (b+a)^{\sigma,i}=b^{\sigma,i}+a^{\sigma,i},
  \qquad \sigma\in\mathbb Z[n],\ i\in[n].
\end{equation*}

\VerdierSemanticAnchor{ch1-2-5-3}
\noindent\textbf{2.5.3.}\quad
Soient $f,g\colon X\rightrightarrows Y$ deux morphismes de complexes,
$a\colon f\rightsquigarrow g$ une homotopie et
$u\colon Y\longrightarrow Y'$ (resp. $v\colon X'\longrightarrow X$) un
morphisme de complexes. On désigne par
$u\star a\colon uf\rightsquigarrow ug$ (resp.
$a\star v\colon fv\rightsquigarrow gv$) l'homotopie

\VerdierSemanticAnchor{ch1-2-5-3-whiskering}
\begin{equation*}
  (u\star a)^{\sigma,i}=u^\sigma a^{\sigma,i}
  \qquad
  \text{(resp. }(a\star v)^{\sigma,i}
    =a^{\sigma,i}v^{\sigma+e_i}\text{)}.
\end{equation*}
